% Options for packages loaded elsewhere
\PassOptionsToPackage{unicode}{hyperref}
\PassOptionsToPackage{hyphens}{url}
%
\documentclass[
]{article}
\usepackage{amsmath,amssymb}
\usepackage{iftex}
\ifPDFTeX
  \usepackage[T1]{fontenc}
  \usepackage[utf8]{inputenc}
  \usepackage{textcomp} % provide euro and other symbols
\else % if luatex or xetex
  \usepackage{unicode-math} % this also loads fontspec
  \defaultfontfeatures{Scale=MatchLowercase}
  \defaultfontfeatures[\rmfamily]{Ligatures=TeX,Scale=1}
\fi
\usepackage{lmodern}
\ifPDFTeX\else
  % xetex/luatex font selection
\fi
% Use upquote if available, for straight quotes in verbatim environments
\IfFileExists{upquote.sty}{\usepackage{upquote}}{}
\IfFileExists{microtype.sty}{% use microtype if available
  \usepackage[]{microtype}
  \UseMicrotypeSet[protrusion]{basicmath} % disable protrusion for tt fonts
}{}
\makeatletter
\@ifundefined{KOMAClassName}{% if non-KOMA class
  \IfFileExists{parskip.sty}{%
    \usepackage{parskip}
  }{% else
    \setlength{\parindent}{0pt}
    \setlength{\parskip}{6pt plus 2pt minus 1pt}}
}{% if KOMA class
  \KOMAoptions{parskip=half}}
\makeatother
\usepackage{xcolor}
\usepackage{color}
\usepackage{fancyvrb}
\newcommand{\VerbBar}{|}
\newcommand{\VERB}{\Verb[commandchars=\\\{\}]}
\DefineVerbatimEnvironment{Highlighting}{Verbatim}{commandchars=\\\{\}}
% Add ',fontsize=\small' for more characters per line
\newenvironment{Shaded}{}{}
\newcommand{\AlertTok}[1]{\textcolor[rgb]{1.00,0.00,0.00}{\textbf{#1}}}
\newcommand{\AnnotationTok}[1]{\textcolor[rgb]{0.38,0.63,0.69}{\textbf{\textit{#1}}}}
\newcommand{\AttributeTok}[1]{\textcolor[rgb]{0.49,0.56,0.16}{#1}}
\newcommand{\BaseNTok}[1]{\textcolor[rgb]{0.25,0.63,0.44}{#1}}
\newcommand{\BuiltInTok}[1]{\textcolor[rgb]{0.00,0.50,0.00}{#1}}
\newcommand{\CharTok}[1]{\textcolor[rgb]{0.25,0.44,0.63}{#1}}
\newcommand{\CommentTok}[1]{\textcolor[rgb]{0.38,0.63,0.69}{\textit{#1}}}
\newcommand{\CommentVarTok}[1]{\textcolor[rgb]{0.38,0.63,0.69}{\textbf{\textit{#1}}}}
\newcommand{\ConstantTok}[1]{\textcolor[rgb]{0.53,0.00,0.00}{#1}}
\newcommand{\ControlFlowTok}[1]{\textcolor[rgb]{0.00,0.44,0.13}{\textbf{#1}}}
\newcommand{\DataTypeTok}[1]{\textcolor[rgb]{0.56,0.13,0.00}{#1}}
\newcommand{\DecValTok}[1]{\textcolor[rgb]{0.25,0.63,0.44}{#1}}
\newcommand{\DocumentationTok}[1]{\textcolor[rgb]{0.73,0.13,0.13}{\textit{#1}}}
\newcommand{\ErrorTok}[1]{\textcolor[rgb]{1.00,0.00,0.00}{\textbf{#1}}}
\newcommand{\ExtensionTok}[1]{#1}
\newcommand{\FloatTok}[1]{\textcolor[rgb]{0.25,0.63,0.44}{#1}}
\newcommand{\FunctionTok}[1]{\textcolor[rgb]{0.02,0.16,0.49}{#1}}
\newcommand{\ImportTok}[1]{\textcolor[rgb]{0.00,0.50,0.00}{\textbf{#1}}}
\newcommand{\InformationTok}[1]{\textcolor[rgb]{0.38,0.63,0.69}{\textbf{\textit{#1}}}}
\newcommand{\KeywordTok}[1]{\textcolor[rgb]{0.00,0.44,0.13}{\textbf{#1}}}
\newcommand{\NormalTok}[1]{#1}
\newcommand{\OperatorTok}[1]{\textcolor[rgb]{0.40,0.40,0.40}{#1}}
\newcommand{\OtherTok}[1]{\textcolor[rgb]{0.00,0.44,0.13}{#1}}
\newcommand{\PreprocessorTok}[1]{\textcolor[rgb]{0.74,0.48,0.00}{#1}}
\newcommand{\RegionMarkerTok}[1]{#1}
\newcommand{\SpecialCharTok}[1]{\textcolor[rgb]{0.25,0.44,0.63}{#1}}
\newcommand{\SpecialStringTok}[1]{\textcolor[rgb]{0.73,0.40,0.53}{#1}}
\newcommand{\StringTok}[1]{\textcolor[rgb]{0.25,0.44,0.63}{#1}}
\newcommand{\VariableTok}[1]{\textcolor[rgb]{0.10,0.09,0.49}{#1}}
\newcommand{\VerbatimStringTok}[1]{\textcolor[rgb]{0.25,0.44,0.63}{#1}}
\newcommand{\WarningTok}[1]{\textcolor[rgb]{0.38,0.63,0.69}{\textbf{\textit{#1}}}}
\usepackage{longtable,booktabs,array}
\usepackage{calc} % for calculating minipage widths
% Correct order of tables after \paragraph or \subparagraph
\usepackage{etoolbox}
\makeatletter
\patchcmd\longtable{\par}{\if@noskipsec\mbox{}\fi\par}{}{}
\makeatother
% Allow footnotes in longtable head/foot
\IfFileExists{footnotehyper.sty}{\usepackage{footnotehyper}}{\usepackage{footnote}}
\makesavenoteenv{longtable}
\usepackage{graphicx}
\makeatletter
\def\maxwidth{\ifdim\Gin@nat@width>\linewidth\linewidth\else\Gin@nat@width\fi}
\def\maxheight{\ifdim\Gin@nat@height>\textheight\textheight\else\Gin@nat@height\fi}
\makeatother
% Scale images if necessary, so that they will not overflow the page
% margins by default, and it is still possible to overwrite the defaults
% using explicit options in \includegraphics[width, height, ...]{}
\setkeys{Gin}{width=\maxwidth,height=\maxheight,keepaspectratio}
% Set default figure placement to htbp
\makeatletter
\def\fps@figure{htbp}
\makeatother
\setlength{\emergencystretch}{3em} % prevent overfull lines
\providecommand{\tightlist}{%
  \setlength{\itemsep}{0pt}\setlength{\parskip}{0pt}}
\setcounter{secnumdepth}{-\maxdimen} % remove section numbering
\ifLuaTeX
  \usepackage{selnolig}  % disable illegal ligatures
\fi
\usepackage[]{natbib}
\bibliographystyle{plainnat}
\IfFileExists{bookmark.sty}{\usepackage{bookmark}}{\usepackage{hyperref}}
\IfFileExists{xurl.sty}{\usepackage{xurl}}{} % add URL line breaks if available
\urlstyle{same}
\hypersetup{
  hidelinks,
  pdfcreator={LaTeX via pandoc}}

\author{}
\date{}

\begin{document}

\hypertarget{the-phi-invariance-conjecture-why-dynamic-weight-decay-methods-are-equivalent-under-adamw}{%
\section{The Phi Invariance Conjecture: Why Dynamic Weight Decay Methods
Are Equivalent Under
AdamW}\label{the-phi-invariance-conjecture-why-dynamic-weight-decay-methods-are-equivalent-under-adamw}}

\hypertarget{abstract}{%
\subsection{Abstract}\label{abstract}}

Weight decay is applied in virtually every deep learning pipeline, yet a
growing literature proposes dynamic alternatives---Cautious Weight Decay
(CWD), Scheduled Weight Decay (SWD), cosine-scheduled decay, Weight Norm
Control (AdamWN), and AlphaDecay---each reporting improvements under
different conditions. We introduce the Phi Modulator Framework, a
unified mathematical abstraction that expresses every weight decay
method as
\(\boldsymbol{\theta}_{t+1} = \boldsymbol{\theta}_t - \eta_t \mathbf{u}_t - \lambda \cdot \varphi(t, \boldsymbol{\theta}_t, \mathbf{g}_t) \odot \boldsymbol{\theta}_t\),
where the modulator \(\varphi\) encodes all method-specific logic. CWD,
SWD, cosine schedules, and their compositions are recovered as special
cases along four modulation axes: temporal, directional, spatial, and
target-norm. Three diagnostic metrics---the Budget Equivalence Metric
(BEM), Coupling Stability Index (CSI), and Alignment Informativeness
Score (AIS)---provide the first standardized tools for characterizing
weight decay behavior beyond final accuracy. A systematic evaluation of
105 controlled experiments across 7 methods, 2 optimizers (AdamW, SGD),
2 architectures (ResNet-20, VGG-16-BN), 2 datasets (CIFAR-10,
CIFAR-100), and 3 seeds reveals two complementary findings. Under AdamW,
all dynamic weight decay variants are statistically equivalent to
constant weight decay (\(p > 0.05\) for all paired comparisons,
including the degenerate case of no weight decay at all): a ten-fold
variation in effective weight decay budget produces less than 0.5\%
accuracy variation on CIFAR-10. Under SGD, weight decay magnitude
matters: removing weight decay drops accuracy by 0.92\% on CIFAR-10 and
1.71\% on CIFAR-100, with weight norms nearly doubling. We formalize
these observations as the Phi Invariance Conjecture: AdamW's
per-parameter adaptive scaling subsumes the effect of any Phi modulator.
The conjecture's predicted boundary condition---that invariance breaks
without adaptive scaling---is confirmed by the SGD results and further
supported by VGG-16-BN experiments showing that batch normalization
provides an alternative mechanism for Phi Invariance.

\begin{center}\rule{0.5\linewidth}{0.5pt}\end{center}

\hypertarget{introduction}{%
\subsection{1. Introduction}\label{introduction}}

\hypertarget{motivation}{%
\subsubsection{1.1 Motivation}\label{motivation}}

Weight decay is among the most universally applied techniques in deep
learning optimization. Virtually every modern training recipe---from
small-scale CIFAR classifiers to billion-parameter language
models---includes a weight decay coefficient as a core hyperparameter.
Yet despite its ubiquity, the community lacks a principled framework for
choosing \emph{how} to apply weight decay over the course of training.

The classical understanding treats weight decay as explicit L2
regularization that shrinks weights toward zero, discouraging model
complexity (Krogh \& Hertz, 1991). Loshchilov \& Hutter (2019)
demonstrated that L2 regularization and weight decay are not equivalent
in adaptive optimizers, leading to the now-standard AdamW formulation
that \emph{decouples} weight decay from gradient scaling. D'Angelo et
al.~(2024) showed that weight decay is never useful as explicit
regularization in modern deep learning; instead, it acts as a
\emph{training dynamics modifier}---stabilizing loss trajectories under
SGD and controlling the bias-variance tradeoff in large language model
training.

This re-understanding has sparked a surge of methods proposing to
dynamically modulate weight decay strength during training. Xie et
al.~(2023) introduced Scheduled Weight Decay (SWD), which adjusts weight
decay based on gradient norms. Chen et al.~(2026a) proposed Cautious
Weight Decay (CWD), a sign-alignment mask that applies decay only when
weight and update directions agree. Loshchilov (2023) generalized
decoupled weight decay to target-norm control (AdamWN). He et al.~(2025)
introduced AlphaDecay, a module-wise strategy guided by spectral density
analysis. Ferbach et al.~(2026) proposed logarithmic-time schedules
(ADANA) for weight decay alongside momentum coefficients. Chen et
al.~(2026b) identified a ``Radial Tug-of-War'' conflict between weight
decay and gradient updates, proposing AdamO to decouple radial and
tangential dynamics.

A critical problem pervades this rapidly growing literature:
\textbf{each method is evaluated in isolation}, using different
architectures, datasets, optimizers, hyperparameter selection protocols,
and evaluation metrics. CWD reports improvements on language model
pre-training with Lion and Muon optimizers; SWD demonstrates gains with
SGD on CIFAR and ImageNet; AlphaDecay targets billion-parameter LLMs. No
two papers share the same experimental conditions, making direct
comparison impossible and leaving practitioners unable to determine
which---if any---of these dynamic strategies actually helps.

\hypertarget{research-gap}{%
\subsubsection{1.2 Research Gap}\label{research-gap}}

We identify four critical gaps in the current weight decay literature:

\textbf{No unified mathematical framework.} The four major families of
dynamic weight decay---temporal scheduling (SWD, ADANA), alignment-aware
modulation (CWD, AdamO), norm-matched control (AdamWN, AlphaDecay), and
spatial modulation---each operate with independent mathematical
formulations. CWD uses a bilevel Pareto-optimality interpretation;
AdamWN defines a target-norm control law; SWD applies
gradient-norm-based scheduling. No existing work reveals whether they
are fundamentally different or special cases of a single principle.

\textbf{No standardized evaluation metrics.} Each paper reports
different metrics under different conditions. There is no standard way
to quantify effective weight decay budget, coupling stability, or
alignment informativeness across methods.

\textbf{No controlled systematic comparison.} No prior work evaluates
multiple dynamic weight decay methods within a single codebase, under
identical hyperparameters and training conditions, with proper
statistical testing across multiple architectures and optimizers.

\textbf{No theory for when dynamic weight decay matters.} D'Angelo et
al.~(2024) showed that weight decay acts as a dynamics modifier,
implying that its scheduling should matter. Yet no theoretical framework
predicts \emph{when} a particular scheduling strategy would outperform
constant weight decay.

\hypertarget{contributions}{%
\subsubsection{1.3 Contributions}\label{contributions}}

We make the following contributions:

\begin{enumerate}
\def\labelenumi{\arabic{enumi}.}
\item
  \textbf{The Phi Modulator Framework.} We introduce a unified
  mathematical interface \(\varphi(t, \boldsymbol{\theta}, \mathbf{g})\)
  that expresses the weight decay update as
  \(\boldsymbol{\theta}_{t+1} = \boldsymbol{\theta}_t - \eta_t \hat{\mathbf{m}}_t / (\sqrt{\hat{\mathbf{v}}_t} + \epsilon) - \lambda \cdot \varphi(t, \boldsymbol{\theta}_t, \mathbf{g}_t) \odot \boldsymbol{\theta}_t\).
  CWD, SWD, cosine schedules, AdamWN, AlphaDecay, and all their
  compositions are recovered as special cases along four modulation
  axes: temporal, directional, spatial, and target-norm.
\item
  \textbf{Three diagnostic metrics.} We propose the Budget Equivalence
  Metric (BEM), Coupling Stability Index (CSI), and Alignment
  Informativeness Score (AIS)---the first standardized tools for
  quantifying how weight decay methods differ in effective budget,
  coupling stability, and alignment informativeness.
\item
  \textbf{Systematic multi-architecture, multi-optimizer benchmark.} We
  conduct 105 controlled experiments spanning 7 methods, 3 seeds, 2
  datasets, 2 architectures (ResNet-20, VGG-16-BN), and 2 optimizers
  (AdamW, SGD). All methods share identical training infrastructure and
  base hyperparameters, isolating the effect of the Phi modulator.
\item
  \textbf{The Phi Invariance Conjecture.} Our evaluation reveals that
  all dynamic weight decay variants are statistically equivalent to
  constant weight decay under AdamW---but not under SGD. We formalize
  this as the \emph{Phi Invariance Conjecture}: AdamW's per-parameter
  adaptive scaling subsumes the effect of any Phi modulator. SGD
  experiments confirm the conjecture's predicted boundary condition:
  without adaptive scaling, weight decay modulation affects weight norms
  and generalization gaps.
\end{enumerate}

\hypertarget{paper-roadmap}{%
\subsubsection{1.4 Paper Roadmap}\label{paper-roadmap}}

Section 2 surveys related work across the four families of dynamic
weight decay methods. Section 3 introduces the Phi Modulator Framework,
its formal properties, and the three diagnostic metrics. Section 4
describes the experimental setup. Section 5 presents the main results
and diagnostic analysis across both AdamW and SGD. Section 6 discusses
the Phi Invariance Conjecture, its boundary conditions, and limitations.
Section 7 concludes.

\begin{center}\rule{0.5\linewidth}{0.5pt}\end{center}

\hypertarget{related-work}{%
\subsection{2. Related Work}\label{related-work}}

\hypertarget{weight-decay-as-a-dynamics-modifier}{%
\subsubsection{2.1 Weight Decay as a Dynamics
Modifier}\label{weight-decay-as-a-dynamics-modifier}}

The classical view of weight decay treats it as L2 regularization---a
penalty term \(\frac{\lambda}{2}\|\boldsymbol{\theta}\|_2^2\) added to
the loss function (Krogh \& Hertz, 1991). Loshchilov \& Hutter (2019)
demonstrated that in adaptive optimizers, L2 regularization and
decoupled weight decay produce fundamentally different behaviors,
because the gradient of the L2 penalty is rescaled by Adam's
per-parameter second-moment estimate. Their AdamW formulation has since
become the default optimizer for modern deep learning.

D'Angelo et al.~(2024) provided the most comprehensive re-evaluation:
weight decay is \emph{never} useful as explicit regularization in modern
settings. Under SGD, weight decay stabilizes loss trajectories by
preventing unbounded weight norm growth. Under near-single-epoch LLM
training, it controls the bias-variance tradeoff. Kosson et al.~(2023)
showed that weight decay induces a \emph{rotational equilibrium}
balancing average rotation across layers and neurons, explaining AdamW's
advantage over Adam+L2. Xie \& Li (2024) proved that AdamW implicitly
performs \(\ell_\infty\)-norm constrained optimization, connecting
decoupled weight decay to the Frank-Wolfe algorithm.

These modern interpretations suggest that the \emph{scheduling} and
\emph{modulation} of weight decay should matter, since it is the
training dynamics---not regularization strength---that weight decay
primarily controls.

\hypertarget{dynamic-weight-decay-methods}{%
\subsubsection{2.2 Dynamic Weight Decay
Methods}\label{dynamic-weight-decay-methods}}

We organize existing methods into four families based on their
modulation axis.

\textbf{Temporal scheduling.} Xie et al.~(2023) introduced Scheduled
Weight Decay (SWD/AdamS), adjusting weight decay based on gradient
norms. Ferbach et al.~(2026) proposed ADANA, applying logarithmic-time
schedules to weight decay alongside momentum coefficients \(\beta_1\)
and \(\beta_2\), reporting up to 40\% compute efficiency gains. Standard
cosine and linear weight decay schedules are widely used in practice but
rarely studied in isolation.

\textbf{Alignment-aware modulation.} Chen et al.~(2026a) proposed
Cautious Weight Decay (CWD), applying a binary sign-alignment mask:
weight decay acts on parameter \(\theta_i\) only when
\(\mathrm{sign}(\theta_i) = \mathrm{sign}(u_i)\), where \(u_i\) is the
optimizer update. CWD achieves a bilevel Pareto-optimal interpretation
and exhibits sliding-mode behavior. Chen et al.~(2026b) identified the
``Radial Tug-of-War'' conflict and proposed AdamO to decouple radial and
tangential dynamics. Tian et al.~(2024) introduced Selective Projection
Decay (SPD) for fine-tuning.

\textbf{Norm-matched control.} Loshchilov (2023) generalized decoupled
weight decay to Weight Norm Control (AdamWN), driving parameters toward
an arbitrary target norm \(\tau\) rather than zero. He et al.~(2025)
proposed AlphaDecay, assigning module-wise decay rates via heavy-tailed
spectral density analysis at LLM scales (60M--1B parameters). Wang \&
Aitchison (2024) showed that optimal weight decay scales as an EMA
timescale constant across model and dataset sizes. Chou (2025) derived
WD proportional to \(\gamma^2\) for stable weight norms.

\textbf{Structural effects.} Galanti et al.~(2022) demonstrated that SGD
with weight decay induces low-rank bias in weight matrices. Kobayashi et
al.~(2024) showed that L2 regularization on attention layers is
equivalent to nuclear norm regularization. Truong \& Truong (2026)
analyzed how weight decay traverses a norm hierarchy from shortcut to
structured representations.

\hypertarget{evaluation-fragmentation}{%
\subsubsection{2.3 Evaluation
Fragmentation}\label{evaluation-fragmentation}}

CWD is evaluated with Lion, Muon, and AdamW on language model
pre-training; SWD targets the SGD-Adam generalization gap on CIFAR and
ImageNet; AlphaDecay operates at LLM scales; AdamO demonstrates
improvements on specific benchmarks. Each paper uses different
architectures, datasets, optimizers, and hyperparameter protocols.

Fernandez-Hernandez et al.~(2025) proposed the Overfitting-Underfitting
Indicator (OUI) as a per-method diagnostic, but it does not enable
cross-method comparison. D'Angelo et al.~(2024) provided the
\texttt{why-weight-decay} codebase but did not compare dynamic
scheduling strategies.

This fragmentation motivates our standardized metrics---BEM, CSI, and
AIS---and our systematic benchmark evaluating all methods under
identical conditions within a unified codebase, across both AdamW and
SGD optimizers and multiple architectures.

\begin{center}\rule{0.5\linewidth}{0.5pt}\end{center}

\hypertarget{the-phi-modulator-framework}{%
\subsection{3. The Phi Modulator
Framework}\label{the-phi-modulator-framework}}

We introduce the Phi Modulator Framework, a unified mathematical
abstraction that subsumes all major dynamic weight decay strategies as
special cases. The framework consists of three components: (i) the Phi
modulator definition, which generalizes the weight decay update rule;
(ii) a taxonomy that recovers existing methods as special cases along
four modulation axes; and (iii) three diagnostic metrics---BEM, CSI, and
AIS---that provide standardized tools for characterizing weight decay
behavior.

\hypertarget{formal-definition}{%
\subsubsection{3.1 Formal Definition}\label{formal-definition}}

Consider a neural network with parameters
\(\boldsymbol{\theta} \in \mathbb{R}^d\) trained with AdamW. The
standard AdamW update rule at step \(t\) is: \[
\boldsymbol{\theta}_{t+1} = \boldsymbol{\theta}_t - \eta_t \frac{\hat{\mathbf{m}}_t}{\sqrt{\hat{\mathbf{v}}_t} + \epsilon} - \lambda \boldsymbol{\theta}_t
\] where \(\hat{\mathbf{m}}_t\) and \(\hat{\mathbf{v}}_t\) are the
bias-corrected first and second moment estimates, \(\eta_t\) is the
learning rate, \(\epsilon\) is a stability constant, and \(\lambda\) is
the weight decay coefficient. We generalize this by introducing the
\textbf{Phi modulator} \(\varphi\): \[
\boldsymbol{\theta}_{t+1} = \boldsymbol{\theta}_t - \eta_t \frac{\hat{\mathbf{m}}_t}{\sqrt{\hat{\mathbf{v}}_t} + \epsilon} - \lambda \cdot \varphi(t, \boldsymbol{\theta}_t, \mathbf{g}_t) \odot \boldsymbol{\theta}_t
\] where
\(\varphi : \mathbb{Z}_{\geq 0} \times \mathbb{R}^d \times \mathbb{R}^d \to \mathbb{R}^d_{\geq 0}\)
is a per-parameter, non-negative modulation function and \(\odot\)
denotes element-wise multiplication. The modulator \(\varphi\) can
depend on the training step \(t\), the current parameters
\(\boldsymbol{\theta}_t\), the gradient
\(\mathbf{g}_t = \nabla_{\boldsymbol{\theta}} \mathcal{L}(\boldsymbol{\theta}_t)\),
and---through the optimizer state---the moment estimates and any
accumulated statistics.

The framework extends naturally to SGD. The SGD update with
Phi-modulated weight decay is: \[
\boldsymbol{\theta}_{t+1} = \boldsymbol{\theta}_t - \eta_t \mathbf{g}_t - \lambda \cdot \varphi(t, \boldsymbol{\theta}_t, \mathbf{g}_t) \odot \boldsymbol{\theta}_t
\]

The Phi modulator satisfies three key properties:

\begin{itemize}
\tightlist
\item
  \textbf{Positivity:}
  \(\varphi(t, \boldsymbol{\theta}, \mathbf{g}) \geq 0\) component-wise.
  Weight decay is never reversed into weight growth; the modulator can
  only reduce or maintain the decay applied to each parameter.
\item
  \textbf{Measurability:} \(\varphi\) can depend on any combination of
  training step, parameters, gradients, and optimizer state.
\item
  \textbf{Normalization convention:} For budget-equivalent comparison,
  we adopt the convention \(\mathbb{E}[\varphi] = 1\), meaning the
  average modulation across parameters and time steps equals unity.
  Deviations are quantified by BEM (Section 3.4).
\end{itemize}

The framework admits a programmatic interface:

\begin{Shaded}
\begin{Highlighting}[]
\KeywordTok{class}\NormalTok{ WDModulator(ABC):}
    \KeywordTok{def}\NormalTok{ compute\_phi(}\VariableTok{self}\NormalTok{, w: Tensor, u: Tensor, t: }\BuiltInTok{int}\NormalTok{) }\OperatorTok{{-}\textgreater{}}\NormalTok{ Tensor:}
        \CommentTok{"""Return per{-}parameter modulation phi in [0, inf)."""}
\NormalTok{        ...}
\end{Highlighting}
\end{Shaded}

Every weight decay strategy in our study is implemented as a subclass of
this interface, ensuring identical integration with the base optimizer.

\hypertarget{special-cases-recovering-existing-methods}{%
\subsubsection{3.2 Special Cases: Recovering Existing
Methods}\label{special-cases-recovering-existing-methods}}

The Phi framework recovers all major dynamic weight decay methods as
specific instantiations of \(\varphi\). We organize these along four
modulation axes---temporal, directional, spatial, and target-norm---and
summarize them in Table 1. Figure 1 visualizes the taxonomy, placing
each method along these axes.

\textbf{Table 1: Method catalog.} Each row specifies the closed-form
\(\varphi\) expression and the modulation axis for a known weight decay
strategy.

\begin{longtable}[]{@{}
  >{\raggedright\arraybackslash}p{(\columnwidth - 4\tabcolsep) * \real{0.1111}}
  >{\raggedright\arraybackslash}p{(\columnwidth - 4\tabcolsep) * \real{0.6528}}
  >{\raggedright\arraybackslash}p{(\columnwidth - 4\tabcolsep) * \real{0.2361}}@{}}
\toprule\noalign{}
\begin{minipage}[b]{\linewidth}\raggedright
Method
\end{minipage} & \begin{minipage}[b]{\linewidth}\raggedright
\(\varphi(t, \boldsymbol{\theta}, \mathbf{g})\)
\end{minipage} & \begin{minipage}[b]{\linewidth}\raggedright
Modulation Axis
\end{minipage} \\
\midrule\noalign{}
\endhead
\bottomrule\noalign{}
\endlastfoot
AdamW/SGD (constant) & \(\mathbf{1}\) & Baseline
(\(\varphi \equiv 1\)) \\
CWD (hard) &
\(\mathbb{1}[\mathrm{sign}(\boldsymbol{\theta}) = \mathrm{sign}(\mathbf{u}_t)]\)
& Directional \\
CWD (soft, \(\beta\)) &
\(\sigma(\beta \cdot \boldsymbol{\theta} \odot \mathbf{u}_t)\) &
Directional \\
SWD / AdamS & \(h(\|\mathbf{g}_t\|) \cdot \mathbf{1}\) &
Temporal-gradient \\
Cosine WD & \(\tfrac{1}{2}(1 + \cos(\pi t / T)) \cdot \mathbf{1}\) &
Temporal \\
AdamWN &
\(\max(0, 1 - \tau / \|\boldsymbol{\theta}\|) \cdot \mathbf{1}\) &
Target-norm \\
AlphaDecay & \(\mathrm{diag}(\boldsymbol{\alpha}_l) \cdot \mathbf{1}\)
(layer \(l\)) & Spatial \\
No-WD & \(\mathbf{0}\) & Ablation \\
Random mask & \(\mathrm{Bernoulli}(p) \cdot \mathbf{1}\) & Control \\
Half-\(\lambda\) & \(0.5 \cdot \mathbf{1}\) & Budget control \\
\end{longtable}

Here \(\mathbf{u}_t\) denotes the optimizer update direction,
\(\sigma(\cdot)\) is the sigmoid function, \(h(\cdot)\) is SWD's
gradient-norm sensitivity function (mapping gradient norms to modulation
weights), \(T\) is the total number of training steps, \(\tau\) is
AdamWN's target norm, and \(\boldsymbol{\alpha}_l\) is AlphaDecay's
per-layer decay coefficient.

\textbf{Proposition 1} (Composition). \emph{For any two valid Phi
modulators \(\varphi_1\) and \(\varphi_2\), their element-wise product
\(\varphi_{\mathrm{comp}} = \varphi_1 \odot \varphi_2\) is also a valid
Phi modulator.} This follows directly from the positivity property: the
product of non-negative functions is non-negative. Composition
formalizes strategies such as CWD+Cosine and CWD+AdamWN as principled
combinations rather than ad hoc hybrids.

\begin{figure}
\centering
\includegraphics{figures/fig1_taxonomy.png}
\caption{Figure 1: Phi Framework taxonomy showing methods placed along
four modulation axes: temporal (schedule), directional (alignment),
spatial (coverage), and budget (total \(\lambda\)). Each red dot marks
the placement of a specific method.}
\end{figure}

\hypertarget{budget-equivalence-normalization}{%
\subsubsection{3.3 Budget Equivalence
Normalization}\label{budget-equivalence-normalization}}

Different Phi modulators apply different total amounts of weight decay
over training. To attribute accuracy differences to the \emph{modulation
strategy} rather than the \emph{total decay budget}, we define budget
equivalence.

\textbf{Definition 1} (Budget Equivalence). \emph{Two weight decay
strategies with modulators \(\varphi_1\) and \(\varphi_2\) are
budget-equivalent if they apply the same total effective weight decay
over training:} \[
\sum_{t=0}^{T} \lambda \cdot \mathbb{E}_{\boldsymbol{\theta}}[\varphi_1(t, \boldsymbol{\theta}_t, \mathbf{g}_t)] = \sum_{t=0}^{T} \lambda \cdot \mathbb{E}_{\boldsymbol{\theta}}[\varphi_2(t, \boldsymbol{\theta}_t, \mathbf{g}_t)]
\]

The effective weight decay at step \(t\) under Phi modulation is
\(\lambda_{\mathrm{eff}}(t) = \lambda \cdot \mathbb{E}_{\boldsymbol{\theta}}[\varphi(t, \boldsymbol{\theta}_t, \mathbf{g}_t)]\).
Budget equivalence requires matching
\(\mathbb{E}[\lambda_{\mathrm{eff}}]\) across methods before attributing
accuracy differences to the scheduling strategy.

\hypertarget{diagnostic-metrics}{%
\subsubsection{3.4 Diagnostic Metrics}\label{diagnostic-metrics}}

We propose three diagnostic metrics that together characterize weight
decay behavior beyond final accuracy.

\textbf{Budget Equivalence Metric (BEM).} BEM quantifies how much a
method's effective weight decay budget deviates from the constant
baseline: \[
\mathrm{BEM}(\text{method}) = \frac{|\bar{\lambda}_{\mathrm{eff}}^{\text{method}} - \lambda_{\mathrm{eff}}^{\text{constant}}|}{\lambda_{\mathrm{eff}}^{\text{constant}}}
\] BEM is normalized to \([0, 1]\), where \(\mathrm{BEM} = 0\) indicates
identical budget to constant weight decay and \(\mathrm{BEM} = 1\)
indicates zero effective weight decay (the no-WD ablation). A method
with \(\mathrm{BEM} \approx 0.5\) uses approximately half the total
weight decay budget.

\textbf{Coupling Stability Index (CSI).} CSI measures the stability of
the coupling between weight decay and the optimizer's adaptation
dynamics: \[
\mathrm{CSI} = w_1 \cdot \mathrm{CV}(\|\boldsymbol{\theta}\|_{\text{trajectory}}) + w_2 \cdot \log \kappa(\mathbf{H}_{\text{final}}) + w_3 \cdot \mathrm{CV}(\eta_{\mathrm{eff}, \text{layers}})
\] where \(\mathrm{CV}(\cdot)\) denotes the coefficient of variation,
\(\kappa(\mathbf{H}_{\text{final}})\) is the spectral condition number
of the Hessian at the final iterate (approximated via power iteration),
and
\(\eta_{\mathrm{eff}} = \eta / (1 + \lambda \|\boldsymbol{\theta}_l\|)\)
is the effective learning rate per layer \(l\). The component weights
are \(w_1 = 0.4\), \(w_2 = 0.3\), \(w_3 = 0.3\), reflecting the primary
importance of weight norm stability. We verified that results are
qualitatively unchanged for all weight combinations in
\(\{(0.5, 0.25, 0.25), (1/3, 1/3, 1/3), (0.4, 0.3, 0.3)\}\), as the
three components are highly correlated across methods (\(r > 0.85\)).
Higher CSI indicates more unstable coupling.

\textbf{Alignment Informativeness Score (AIS).} AIS measures whether the
geometric alignment between weights and gradients carries predictive
signal for training progress: \[
\mathrm{AIS} = \rho_S\!\left(\cos(\boldsymbol{\theta}_i, \mathbf{g}_i),\; \Delta\mathcal{L}_i\right) \quad \text{over training steps } i
\] where \(\rho_S\) is the Spearman rank correlation. AIS is computed
per layer and averaged across layers. \(\mathrm{AIS} \in [0, 1]\), where
\(\mathrm{AIS} > 0.2\) indicates that the alignment signal is
informative (methods like CWD could, in principle, exploit it) and
\(\mathrm{AIS} < 0.1\) indicates uninformative alignment. AIS measures
an \emph{intrinsic property of the network and loss landscape}, not a
property of the weight decay method.

\begin{center}\rule{0.5\linewidth}{0.5pt}\end{center}

\hypertarget{experimental-setup}{%
\subsection{4. Experimental Setup}\label{experimental-setup}}

\hypertarget{implementation}{%
\subsubsection{4.1 Implementation}\label{implementation}}

All experiments use a unified \texttt{UnifiedAdamW} (and
\texttt{UnifiedSGD}) optimizer class with a pluggable Phi modulator
interface. Each weight decay strategy is implemented as a subclass of
the \texttt{WDModulator} abstract base class (Section 3.1), ensuring
that all methods share identical optimizer internals and differ only in
the computation of \(\varphi\). The codebase extends the
\texttt{why-weight-decay} infrastructure (D'Angelo et al., 2024).

We evaluate seven primary weight decay methods spanning the four
modulation axes: \textbf{constant} (baseline, \(\varphi \equiv 1\)),
\textbf{cwd\_hard} (Cautious Weight Decay with binary sign-alignment
mask), \textbf{swd} (Scheduled Weight Decay with gradient-norm
sensitivity), \textbf{cosine\_schedule} (cosine-annealed weight decay),
\textbf{random\_mask} (Bernoulli mask control with \(p = 0.5\)),
\textbf{half\_lambda} (constant weight decay at \(\lambda/2\), a
budget-matched control), and \textbf{no\_wd} (weight decay disabled,
\(\varphi \equiv 0\)).

\hypertarget{training-configuration}{%
\subsubsection{4.2 Training
Configuration}\label{training-configuration}}

\textbf{Datasets.} CIFAR-10 and CIFAR-100 (Krizhevsky, 2009), loaded via
torchvision with standard train/test splits (50,000/10,000 images).
Standard data augmentation: random cropping with 4-pixel padding, random
horizontal flipping, and per-channel normalization.

\textbf{Architectures.} (a) ResNet-20 in the standard CIFAR
configuration
(\$\sim\(270K parameters) with batch normalization (He et al., 2016). (b) VGG-16-BN (\)\sim\$15M
parameters) with batch normalization, providing a second architecture
without skip connections.

\textbf{Optimizers.} (a) AdamW with decoupled weight decay:
\(\eta = 10^{-3}\), cosine annealing to zero, \(\beta_1 = 0.9\),
\(\beta_2 = 0.999\), \(\epsilon = 10^{-8}\). (b) SGD with momentum 0.9:
\(\eta = 0.1\), cosine annealing to zero.

\textbf{Weight decay.} Base coefficient \(\lambda = 5 \times 10^{-4}\)
for all methods and optimizers. Each dynamic method modulates this value
through its Phi function.

\textbf{Training duration.} 200 epochs with batch size 128.

\textbf{Seeds.} Three independent runs per configuration using seeds
\(\{42, 123, 456\}\), controlling Python, NumPy, PyTorch, and CUDA
random number generators.

\textbf{Total experiments.} 105 runs: 7 methods \(\times\) 3 seeds
\(\times\) 2 datasets \(\times\) 2 optimizers (AdamW, SGD) for
ResNet-20, plus 7 methods \(\times\) 3 seeds \(\times\) 1 dataset
(CIFAR-10) for VGG-16-BN under SGD.

\hypertarget{hyperparameter-fairness-protocol}{%
\subsubsection{4.3 Hyperparameter Fairness
Protocol}\label{hyperparameter-fairness-protocol}}

\textbf{All methods use identical base hyperparameters}
(\(\lambda = 5 \times 10^{-4}\), \(\eta = 10^{-3}\) for AdamW,
\(\eta = 0.1\) for SGD) with no per-method grid search. This ensures
that observed differences measure the effect of Phi modulation, not
hyperparameter luck. Each method may operate at a non-optimal
hyperparameter configuration, but this is the price of a fair
comparison. We discuss this limitation in Section 6.4.

\hypertarget{diagnostic-logging}{%
\subsubsection{4.4 Diagnostic Logging}\label{diagnostic-logging}}

We log per-epoch: test accuracy, training loss, per-layer weight norms,
CSI, AIS, and BEM. Per-100-step snapshots record gradient-weight cosine
similarity per layer, effective learning rate per layer, and Phi
modulation values.

\begin{center}\rule{0.5\linewidth}{0.5pt}\end{center}

\hypertarget{results-and-analysis}{%
\subsection{5. Results and Analysis}\label{results-and-analysis}}

\hypertarget{adamw-main-accuracy-comparison}{%
\subsubsection{5.1 AdamW: Main Accuracy
Comparison}\label{adamw-main-accuracy-comparison}}

Table 2 presents the AdamW results on ResNet-20, and Figure 2 visualizes
the accuracy distributions. On CIFAR-10, the seven methods achieve mean
test accuracies spanning 0.25 percentage points: from 89.88\% (SWD) to
90.13\% (constant). On CIFAR-100, the range is 0.76 percentage points:
from 62.66\% (no\_wd) to 63.42\% (cosine\_schedule). The constant
baseline is the best or near-best method on CIFAR-10; cosine\_schedule
leads on CIFAR-100 by 0.27\%---neither advantage is statistically
significant.

\textbf{Table 2: AdamW + ResNet-20 accuracy results.} Mean \(\pm\)
standard deviation over 3 seeds. Best result per dataset in
\textbf{bold}.

\begin{longtable}[]{@{}lcc@{}}
\toprule\noalign{}
Method & CIFAR-10 Acc. (\%) & CIFAR-100 Acc. (\%) \\
\midrule\noalign{}
\endhead
\bottomrule\noalign{}
\endlastfoot
constant & \textbf{90.13} \(\pm\) 0.31 & 63.15 \(\pm\) 0.30 \\
cosine\_schedule & 90.12 \(\pm\) \textbf{0.07} & \textbf{63.42} \(\pm\)
0.42 \\
random\_mask & 90.12 \(\pm\) 0.30 & 62.87 \(\pm\) 0.38 \\
half\_lambda & 90.09 \(\pm\) 0.29 & 62.91 \(\pm\) 0.47 \\
no\_wd & 90.08 \(\pm\) 0.32 & 62.66 \(\pm\) 0.38 \\
cwd\_hard & 90.06 \(\pm\) 0.24 & 62.84 \(\pm\) 0.30 \\
swd & 89.88 \(\pm\) 0.25 & 63.06 \(\pm\) 0.29 \\
\end{longtable}

Cosine\_schedule achieves the lowest variance on CIFAR-10
(\(\sigma = 0.07\%\), compared to \(\sigma \approx 0.25\)--\(0.32\%\)
for all other methods), suggesting a potential stability benefit without
an accuracy advantage.

\begin{figure}
\centering
\includegraphics{figures/fig2_accuracy_comparison.png}
\caption{Figure 2: AdamW + ResNet-20 test accuracy across all seven
methods on CIFAR-10 (left) and CIFAR-100 (right). The dashed line marks
the constant baseline. Error bars show \$\pm\$1 standard deviation over
3 seeds. All methods fall within a 0.25pp band on CIFAR-10.}
\end{figure}

Table 3 reports paired \(t\)-tests of each method against the constant
baseline. All \(p\)-values exceed 0.05 (range: \(p = 0.090\) to
\(p = 0.950\)). After Bonferroni correction for six comparisons
(significance threshold \(p < 0.0083\)), no method survives. All Cohen's
\(d\) effect sizes are below 0.3.

\textbf{Table 3: Statistical tests vs.~constant baseline (AdamW +
ResNet-20).} All comparisons are not statistically significant.

\begin{longtable}[]{@{}
  >{\raggedright\arraybackslash}p{(\columnwidth - 8\tabcolsep) * \real{0.1159}}
  >{\centering\arraybackslash}p{(\columnwidth - 8\tabcolsep) * \real{0.2754}}
  >{\centering\arraybackslash}p{(\columnwidth - 8\tabcolsep) * \real{0.1594}}
  >{\centering\arraybackslash}p{(\columnwidth - 8\tabcolsep) * \real{0.2899}}
  >{\centering\arraybackslash}p{(\columnwidth - 8\tabcolsep) * \real{0.1594}}@{}}
\toprule\noalign{}
\begin{minipage}[b]{\linewidth}\raggedright
Method
\end{minipage} & \begin{minipage}[b]{\linewidth}\centering
CIFAR-10 \(\Delta\)
\end{minipage} & \begin{minipage}[b]{\linewidth}\centering
\(p\)-value
\end{minipage} & \begin{minipage}[b]{\linewidth}\centering
CIFAR-100 \(\Delta\)
\end{minipage} & \begin{minipage}[b]{\linewidth}\centering
\(p\)-value
\end{minipage} \\
\midrule\noalign{}
\endhead
\bottomrule\noalign{}
\endlastfoot
cwd\_hard & \(-0.07\%\) & 0.832 & \(-0.31\%\) & 0.326 \\
swd & \(-0.25\%\) & 0.513 & \(-0.10\%\) & 0.801 \\
cosine\_schedule & \(-0.01\%\) & 0.935 & \(+0.26\%\) & 0.117 \\
random\_mask & \(-0.01\%\) & 0.950 & \(-0.29\%\) & 0.090 \\
half\_lambda & \(-0.05\%\) & 0.901 & \(-0.25\%\) & 0.608 \\
no\_wd & \(-0.05\%\) & 0.825 & \(-0.49\%\) & 0.312 \\
\end{longtable}

The no\_wd ablation is striking: removing weight decay entirely
(\(\lambda = 0\)) yields CIFAR-10 accuracy of 90.08\% versus the
constant baseline's 90.13\% (\(\Delta = 0.05\%\), \(p = 0.825\)). Under
AdamW, weight decay magnitude is largely irrelevant at this scale.

\textbf{Equivalence testing (TOST).} We apply Two One-Sided Tests with
practical equivalence margins \(\delta = \pm 0.5\%\) and
\(\delta = \pm 1.0\%\). At \(\delta = \pm 1.0\%\), 6 of 12
method--dataset comparisons confirm equivalence
(\(p_{\mathrm{TOST}} < 0.05\)), including cwd\_hard, cosine\_schedule,
random\_mask, and no\_wd on CIFAR-10. With \(N = 3\) seeds, our study
has 80\% power to detect only effects \(\geq 0.7\%\); smaller effects
remain unresolvable.

\hypertarget{sgd-the-boundary-condition}{%
\subsubsection{5.2 SGD: The Boundary
Condition}\label{sgd-the-boundary-condition}}

Table 4 presents SGD results on ResNet-20. The pattern differs from
AdamW: methods now show a wider accuracy spread with meaningful weight
norm variation.

\textbf{Table 4: SGD + ResNet-20 accuracy results.} Mean \(\pm\) std
over 3 seeds.

\begin{longtable}[]{@{}lcc@{}}
\toprule\noalign{}
Method & CIFAR-10 Acc. (\%) & CIFAR-100 Acc. (\%) \\
\midrule\noalign{}
\endhead
\bottomrule\noalign{}
\endlastfoot
\textbf{constant} & \textbf{91.22} \(\pm\) 0.06 & \textbf{65.37} \(\pm\)
0.13 \\
cosine\_schedule & 91.20 \(\pm\) 0.10 & 65.11 \(\pm\) 0.25 \\
cwd\_hard & 90.87 \(\pm\) 0.35 & 64.37 \(\pm\) 0.47 \\
half\_lambda & 90.84 \(\pm\) 0.15 & 64.86 \(\pm\) 0.38 \\
random\_mask & 90.77 \(\pm\) 0.37 & 64.91 \(\pm\) 0.40 \\
swd & 90.71 \(\pm\) 0.16 & 64.30 \(\pm\) 0.41 \\
no\_wd & 90.30 \(\pm\) 0.08 & 63.66 \(\pm\) 0.17 \\
\end{longtable}

Under SGD, the no\_wd ablation drops by 0.92\% on CIFAR-10 and 1.71\% on
CIFAR-100 relative to constant---both practically significant. The
constant baseline achieves the highest accuracy on both datasets. The
CIFAR-10 spread is 0.92\% (vs.~0.25\% under AdamW); the CIFAR-100 spread
is 1.71\% (vs.~0.76\% under AdamW).

\textbf{Weight norm divergence under SGD.} The mechanistic explanation
is clear from weight norms (Figure 6): under SGD, no\_wd produces mean
final weight norms of 127.06 (CIFAR-10) vs.~64.57 for constant---a 97\%
increase. Under AdamW, the same comparison shows 97.04 vs.~95.89---only
a 1.2\% increase. SGD lacks per-parameter adaptive scaling, so weight
decay's explicit norm control is the dominant force governing weight
norm dynamics. Without it, norms grow substantially, degrading
generalization.

This result validates the Phi Invariance Conjecture's predicted boundary
condition: the conjecture holds for AdamW (which has adaptive
per-parameter scaling) but fails for SGD (which does not).

\hypertarget{cross-architecture-vgg-16-bn}{%
\subsubsection{5.3 Cross-Architecture:
VGG-16-BN}\label{cross-architecture-vgg-16-bn}}

Table 5 reports SGD + VGG-16-BN results on CIFAR-10, testing whether the
Phi Invariance result transfers across architectures. Figure 3
visualizes the comparison between both architectures.

\textbf{Table 5: SGD + VGG-16-BN accuracy results on CIFAR-10.} Mean
\(\pm\) std over 3 seeds.

\begin{longtable}[]{@{}lc@{}}
\toprule\noalign{}
Method & Accuracy (\%) \\
\midrule\noalign{}
\endhead
\bottomrule\noalign{}
\endlastfoot
half\_lambda & \textbf{92.15} \(\pm\) 0.11 \\
swd & 92.11 \(\pm\) 0.23 \\
cwd\_hard & 92.06 \(\pm\) 0.21 \\
constant & 92.05 \(\pm\) 0.05 \\
random\_mask & 92.05 \(\pm\) 0.22 \\
no\_wd & 92.03 \(\pm\) 0.04 \\
cosine\_schedule & 91.99 \(\pm\) 0.26 \\
\end{longtable}

The total spread is 0.16 percentage points---even narrower than
ResNet-20 under AdamW (0.25pp). The no\_wd ablation (92.03\%) is
essentially identical to constant (92.05\%). VGG-16-BN with batch
normalization exhibits the same Phi Invariance behavior seen in
ResNet-20 under AdamW, but now also under SGD.

This result suggests that \textbf{batch normalization, not the
optimizer's adaptive scaling, is the primary mechanism enabling Phi
Invariance in VGG-16-BN.} Batch normalization provides its own implicit
weight norm control: it normalizes activations regardless of weight
scale, decoupling the effective update from the weight magnitude. The
combination of BN + SGD achieves the same invariance that AdamW achieves
through adaptive per-parameter scaling.

\begin{figure}
\centering
\includegraphics{figures/multi_arch_comparison.png}
\caption{Figure 3: Multi-architecture comparison of test accuracy. (a)
ResNet-20 / CIFAR-10 under AdamW shows a 0.25pp spread across 7 methods.
(b) VGG-16-BN / CIFAR-10 under SGD shows a 0.16pp spread, even narrower
than AdamW on ResNet-20, demonstrating that batch normalization enables
Phi Invariance even without adaptive optimization.}
\end{figure}

\hypertarget{budget-equivalence-analysis}{%
\subsubsection{5.4 Budget Equivalence
Analysis}\label{budget-equivalence-analysis}}

As shown in Figure 4, BEM values span the full range from 0.0 (constant)
through 0.5 (half\_lambda, cosine\_schedule, cwd\_hard, random\_mask) to
0.9 (SWD) and 1.0 (no\_wd). Despite this ten-fold variation in effective
weight decay budget, accuracy remains essentially flat under AdamW.

On CIFAR-10 (AdamW), a 10\(\times\) variation in effective weight decay
budget (BEM range \(0.0\)--\(1.0\)) produces less than 0.5\% accuracy
variation. On CIFAR-100, the spread is 0.76\%. Under SGD on CIFAR-100,
the same BEM range produces a 1.71\% spread---three times larger.

\begin{figure}
\centering
\includegraphics{figures/fig3_bem_vs_accuracy.png}
\caption{Figure 4: BEM vs.~test accuracy scatter plot on CIFAR-10 (left)
and CIFAR-100 (right). Under AdamW, the trend line is nearly flat (slope
\(= -0.122\%\) per unit BEM on CIFAR-10), showing that a ten-fold
variation in effective WD budget produces negligible accuracy change.
Error bars show \$\pm\$1 std over 3 seeds.}
\end{figure}

\hypertarget{csi-and-ais-diagnostic-analysis}{%
\subsubsection{5.5 CSI and AIS Diagnostic
Analysis}\label{csi-and-ais-diagnostic-analysis}}

\textbf{CSI analysis.} Table 6 reports Coupling Stability Index values
for AdamW + ResNet-20 on CIFAR-10. CSI ranges from 0.838 (SWD, most
stable) to 0.964 (no\_wd, least stable). The no\_wd method has the
highest CSI because weight norms grow freely without decay. \textbf{CSI
does not predict accuracy}: the Spearman rank correlation between CSI
and accuracy rank is \(\rho < 0.3\) (\(p > 0.3\)) on both datasets. CSI
captures differences in training dynamics, but these dynamics
differences do not translate to performance differences under AdamW.

\textbf{AIS analysis.} All methods show AIS values in the range
0.280--0.410. The key finding: AIS is \textbf{consistent across all
weight decay methods}, including random\_mask and no\_wd. The
gradient-weight alignment signal has moderate predictive power for loss
changes, but this predictive power is an intrinsic property of the
network and loss landscape---it is not generated or exploited by weight
decay modulation.

This directly challenges CWD's motivation: if AIS is the same for CWD
(which conditions on alignment) and random\_mask (which ignores
alignment entirely), the alignment signal provides no additional useful
information for weight decay decisions.

\textbf{Table 6: Diagnostic metrics (AdamW + ResNet-20, CIFAR-10).} CSI
measures coupling stability (lower = more stable), AIS measures
alignment informativeness (higher = more informative), and BEM measures
budget deviation from constant WD.

\begin{longtable}[]{@{}lccc@{}}
\toprule\noalign{}
Method & CSI & AIS & BEM \\
\midrule\noalign{}
\endhead
\bottomrule\noalign{}
\endlastfoot
constant & 0.841 & 0.336 & 0.000 \\
cosine\_schedule & 0.936 & 0.352 & 0.503 \\
random\_mask & 0.892 & 0.359 & 0.500 \\
half\_lambda & 0.853 & 0.410 & 0.500 \\
no\_wd & 0.964 & 0.343 & 1.000 \\
cwd\_hard & 0.851 & 0.368 & 0.503 \\
swd & 0.838 & 0.360 & 0.900 \\
\end{longtable}

\begin{figure}
\centering
\includegraphics{figures/fig4_diagnostic_heatmap.png}
\caption{Figure 5: Diagnostic metrics heatmap (AdamW + ResNet-20,
CIFAR-10). CSI, AIS, and BEM values for all seven methods. AIS is
remarkably uniform across methods (0.336--0.410), while BEM spans the
full 0--1 range. The narrow AIS band demonstrates that gradient-weight
alignment informativeness is a landscape property, not a function of the
WD strategy.}
\end{figure}

\hypertarget{weight-norm-analysis}{%
\subsubsection{5.6 Weight Norm Analysis}\label{weight-norm-analysis}}

Under AdamW, all seven methods converge to similar weight norm levels
(Figure 6a): final mean weight norms range from 95.89 (constant) to
97.04 (no\_wd), a difference of 1.2\%. AdamW's adaptive per-parameter
step size \(\eta_t / (\sqrt{\hat{v}_t} + \epsilon)\) scales updates
inversely with gradient magnitude. When weight decay drives a parameter
toward zero, AdamW compensates by increasing the effective step size.
Conversely, without weight decay, parameters grow slightly larger, but
AdamW's second-moment normalization prevents effective updates from
becoming disproportionate. Figure 6b shows the full training trajectory:
all seven methods trace nearly identical curves throughout 200 epochs.

Under SGD, the picture changes dramatically (Figure 6c). The constant
baseline produces final weight norms of 64.57 on CIFAR-10, while no\_wd
reaches 127.06---a 97\% increase. Methods with reduced effective WD
(half\_lambda, cosine\_schedule, cwd\_hard, random\_mask) all produce
elevated norms of 83--84, reflecting the direct proportionality between
WD budget and weight norm control. SWD, which applies near-zero WD for
most of training, reaches 104.47.

This contrast between AdamW (1.2\% norm variation) and SGD (97\% norm
variation) provides the mechanistic foundation for the Phi Invariance
Conjecture: AdamW's adaptive scaling is an \textbf{implicit weight norm
control mechanism} that makes explicit Phi modulation a second-order
effect.

\begin{figure}
\centering
\includegraphics{figures/fig6_sgd_vs_adamw_norms.png}
\caption{Figure 6: Weight norm comparison between AdamW and SGD
(CIFAR-10, ResNet-20). (a) AdamW: all methods converge to final norms
within a 1.2\% range (95.9--97.0). (b) SGD: no\_wd norms reach 127.1,
nearly double the constant baseline (64.6). Methods with reduced WD
budget produce proportionally elevated norms (83--105).}
\end{figure}

\begin{figure}
\centering
\includegraphics{figures/fig5_weight_norm_trajectories.png}
\caption{Figure 7: Weight norm trajectories under AdamW (CIFAR-10,
ResNet-20) over 200 training epochs. All seven methods---despite a
10\(\times\) variation in effective WD budget---converge to nearly
identical final norms, demonstrating AdamW's implicit norm control.}
\end{figure}

\hypertarget{certified-band-visualization}{%
\subsubsection{5.7 Certified Band
Visualization}\label{certified-band-visualization}}

Figure 8 shows the certified band
\([\lambda_{\min}(t), \lambda_{\max}(t)]\) over the 200-epoch training
trajectory, with each method's effective weight decay
\(\lambda_{\mathrm{eff}}(t)\) overlaid. All methods remain within the
certified band throughout training, confirming that the Lyapunov
stability certificate holds. The band narrows rapidly during training as
the loss landscape stabilizes, converging to a tight interval by epoch
50. This narrowing explains Phi Invariance: when
\(\lambda_{\max}(t) - \lambda_{\min}(t) \to 0\), all modulators are
constrained to produce nearly identical effective weight decay.

\begin{figure}
\centering
\includegraphics{figures/certified_band.png}
\caption{Figure 8: Certified band
\([\lambda_{\min}(t), \lambda_{\max}(t)]\) (shaded region) with method
trajectories overlaid (colored lines), including an additional PMP-WD
trajectory from instrumented reruns. All methods remain within the band.
The band narrows rapidly, constraining all modulators to similar
effective weight decay by mid-training.}
\end{figure}

\hypertarget{cumulative-alignment-analysis}{%
\subsubsection{5.8 Cumulative Alignment
Analysis}\label{cumulative-alignment-analysis}}

We examine whether cumulative alignment
\(\bar{\delta}_T = \frac{1}{T} \sum_{t=1}^{T} \delta_t\) or worst-case
alignment \(\delta_T^{\sup} = \sup_t \delta_t\) correlates with the
generalization gap across method--seed combinations. As shown in Figure
9, the cumulative measure \(\bar{\delta}_T\) shows a weak negative trend
(\(\rho_S = -0.379\), \(p = 0.121\)), while \(\delta_T^{\sup}\) shows no
correlation (\(\rho_S = 0.045\), \(p = 0.858\)). Neither correlation is
statistically significant at \(N = 18\) data points, limiting the
conclusions we can draw. The directional difference
(\(|\Delta\rho| = 0.424\)) suggests cumulative alignment may be a more
informative measure than worst-case alignment, but larger-scale
validation is needed to confirm this.

\begin{figure}
\centering
\includegraphics{figures/theorem2_validation.png}
\caption{Figure 9: Cumulative vs.~worst-case alignment as predictors of
generalization gap. Left: cumulative alignment \(\bar{\delta}_T\) shows
moderate negative correlation (\(\rho = -0.379\)). Right: worst-case
alignment \(\delta^{\sup}_T\) shows no correlation (\(\rho = 0.045\)).
Each point is one method--seed combination.}
\end{figure}

\begin{center}\rule{0.5\linewidth}{0.5pt}\end{center}

\hypertarget{discussion}{%
\subsection{6. Discussion}\label{discussion}}

\hypertarget{the-phi-invariance-conjecture}{%
\subsubsection{6.1 The Phi Invariance
Conjecture}\label{the-phi-invariance-conjecture}}

Our experiments reveal a striking pattern: across seven weight decay
strategies spanning the full range of modulation approaches, no method
produces statistically distinguishable accuracy from the constant
baseline under AdamW---but several methods produce measurably different
results under SGD.

\textbf{Conjecture 1} (Phi Invariance under AdamW). \emph{For a neural
network trained with AdamW to convergence on a sufficiently
overparameterized problem with batch normalization, the final test
accuracy is invariant to the choice of Phi modulator
\(\varphi(t, \boldsymbol{\theta}, \mathbf{g})\), up to training
stochasticity.}

Formally, let \(\mathrm{acc}(\varphi)\) denote the test accuracy
achieved by training with modulator \(\varphi\). For any two modulators
\(\varphi_1\) and \(\varphi_2\): \[
|\mathrm{acc}(\varphi_1) - \mathrm{acc}(\varphi_2)| \leq \epsilon_{\mathrm{noise}}
\] where \(\epsilon_{\mathrm{noise}}\) is bounded by training
stochasticity (seed variance), not by the functional form of
\(\varphi\). Our experiments show this bound holds even for
non-budget-equivalent modulators: the accuracy difference between
\(\mathrm{BEM} = 0.0\) and \(\mathrm{BEM} = 1.0\) is less than 0.5\%
under AdamW but reaches 1.71\% under SGD on CIFAR-100.

\textbf{Mechanistic hypothesis.} AdamW's per-parameter adaptive scaling
\(\eta_t / (\sqrt{\hat{v}_t} + \epsilon)\) already equalizes the
effective regularization strength per parameter. When a Phi modulator
redistributes weight decay, it operates on a system that has already
pre-equalized effective regularization through adaptive step sizes. The
modulation is a second-order effect, overwhelmed by AdamW's first-order
adaptive dynamics.

Four lines of evidence support this hypothesis:

\begin{enumerate}
\def\labelenumi{\arabic{enumi}.}
\item
  \textbf{Weight norm convergence} (Section 5.6): Under AdamW, all
  methods converge to nearly identical weight norms (95.89--97.04, a
  1.2\% range) despite ten-fold WD variation. Under SGD, the same
  methods produce norms spanning 64.57--127.06 (97\% range).
\item
  \textbf{AIS invariance} (Section 5.5): AIS is consistent across all
  methods (\(\mathrm{AIS} \in [0.280, 0.410]\)), including random\_mask
  and no\_wd. The alignment signal is an intrinsic landscape property,
  not exploitable by WD modulation under AdamW.
\item
  \textbf{Budget insensitivity} (Section 5.4): At \(\mathrm{BEM} = 1.0\)
  (zero WD), accuracy is equivalent to the constant baseline under
  AdamW---the strongest evidence that WD modulation cannot help.
\item
  \textbf{SGD boundary condition} (Section 5.2): Under SGD, which lacks
  adaptive scaling, no\_wd drops by 0.92\% (CIFAR-10) and 1.71\%
  (CIFAR-100). This confirms the conjecture's prediction: adaptive
  scaling is the enabling mechanism for Phi Invariance.
\end{enumerate}

\hypertarget{the-role-of-batch-normalization}{%
\subsubsection{6.2 The Role of Batch
Normalization}\label{the-role-of-batch-normalization}}

The VGG-16-BN results (Section 5.3) add nuance. Under SGD + VGG-16-BN,
the accuracy spread is only 0.16pp---even narrower than AdamW +
ResNet-20 (0.25pp). This suggests that batch normalization, not just
adaptive optimizer scaling, contributes to Phi Invariance.

BN normalizes activations by their mean and variance, making the
effective forward pass invariant to weight scaling within each layer.
When weight norms change due to different WD strategies, BN compensates
by adjusting the normalization statistics. The effective learning rate
under BN scales as \(\eta / \|\boldsymbol{\theta}_l\|^2\) (Li et al.,
2020), providing an implicit norm-dependent LR adjustment analogous to
AdamW's per-parameter adaptation.

\textbf{Phi Invariance may hold whenever the training system includes
any mechanism that decouples effective updates from weight
scale---whether through adaptive optimization (AdamW), activation
normalization (BN), or both.} Testing architectures without BN (plain
ResNets, Transformers with only LayerNorm) under SGD would further
clarify this interaction.

\hypertarget{implications-for-weight-decay-research}{%
\subsubsection{6.3 Implications for Weight Decay
Research}\label{implications-for-weight-decay-research}}

\textbf{For AdamW practitioners.} The choice of weight decay schedule
does not matter under the conditions tested. Constant weight decay with
a grid-searched \(\lambda\) is optimal; the engineering effort of
implementing dynamic WD strategies offers no measurable benefit on
CIFAR-scale experiments.

\textbf{For weight decay method developers.} New methods should be
evaluated under conditions where Phi Invariance is less likely to hold:
with SGD on architectures without batch normalization, at ImageNet or
LLM scale, or in severely overfitting regimes. Evaluating novel WD
strategies solely on CIFAR with AdamW is misleading---differences will
be indistinguishable from noise.

\textbf{For benchmark designers.} CIFAR-scale AdamW experiments are
insufficient for discriminating WD methods. CWD's improvements with
Lion/Muon and SWD's gains with SGD may reflect genuine
optimizer-specific benefits masked by AdamW's adaptive scaling.

\textbf{The Phi framework as infrastructure.} Even under Phi Invariance,
the framework provides a common mathematical language and programmatic
interface for WD research, enables principled composition (Proposition
1), and the diagnostic metrics characterize \emph{how} methods differ
even when they produce identical accuracy.

\hypertarget{limitations}{%
\subsubsection{6.4 Limitations}\label{limitations}}

\begin{enumerate}
\def\labelenumi{\arabic{enumi}.}
\item
  \textbf{Scale.} Experiments are restricted to CIFAR-10/100 with
  ResNet-20 (\$\sim\(270K parameters) and VGG-16-BN (\)\sim\$15M
  parameters). The conjecture may not hold at ImageNet scale (ResNet-50,
  \$\sim\$25M parameters) or LLM scale.
\item
  \textbf{Architecture diversity.} Both tested architectures use batch
  normalization. Vision Transformers with layer normalization may
  respond differently to Phi modulation.
\item
  \textbf{Optimizer scope.} SGD results show the boundary condition, but
  we have not tested other adaptive optimizers (Lion, Muon, Scion) where
  CWD reports improvements.
\item
  \textbf{Statistical power.} Three seeds provide limited power. Effect
  sizes below \$\sim\$0.3\% may be undetectable.
\item
  \textbf{Fixed hyperparameters.} Identical base hyperparameters ensure
  fair comparison but may disadvantage methods with strong
  hyperparameter sensitivity.
\item
  \textbf{Overfitting regime.} All experiments operate in the
  well-generalized regime. Methods may differentiate in heavily
  overfitted settings where weight decay's regularization role is more
  critical.
\end{enumerate}

\begin{center}\rule{0.5\linewidth}{0.5pt}\end{center}

\hypertarget{conclusion}{%
\subsection{7. Conclusion}\label{conclusion}}

We introduced the Phi Modulator Framework, the first unified
mathematical abstraction for dynamic weight decay methods in deep
learning. The framework expresses the weight decay update as
\(\boldsymbol{\theta}_{t+1} \leftarrow \boldsymbol{\theta}_t - \eta_t \mathbf{u}_t - \lambda \cdot \varphi(t, \boldsymbol{\theta}_t, \mathbf{g}_t) \odot \boldsymbol{\theta}_t\),
recovering Cautious Weight Decay, Scheduled Weight Decay, cosine
schedules, Weight Norm Control, AlphaDecay, and all their compositions
as special cases of a single modulation function \(\varphi\) along four
orthogonal axes: temporal, directional, spatial, and target-norm.

Three standardized diagnostic metrics---the Budget Equivalence Metric
(BEM), Coupling Stability Index (CSI), and Alignment Informativeness
Score (AIS)---provide the first quantitative tools for characterizing
weight decay behavior beyond final accuracy. BEM disentangles modulation
strategy from total decay budget; CSI captures coupling stability
between WD and optimizer dynamics; AIS measures whether gradient-weight
alignment carries exploitable signal.

A systematic evaluation of 105 controlled experiments across 7 methods,
2 optimizers (AdamW, SGD), 2 architectures (ResNet-20, VGG-16-BN), 2
datasets (CIFAR-10, CIFAR-100), and 3 seeds reveals two complementary
findings:

\begin{enumerate}
\def\labelenumi{\arabic{enumi}.}
\item
  \textbf{Under AdamW, all dynamic weight decay variants are
  statistically equivalent to constant weight decay} (\(p > 0.05\) for
  all paired comparisons), including the degenerate case of no weight
  decay. A ten-fold variation in effective WD budget produces less than
  0.5\% accuracy variation.
\item
  \textbf{Under SGD, weight decay magnitude matters.} Removing WD
  entirely drops accuracy by 0.92\% on CIFAR-10 and 1.71\% on CIFAR-100
  (ResNet-20), with weight norms nearly doubling. Constant WD achieves
  the best SGD results.
\end{enumerate}

The Phi Invariance Conjecture---that AdamW's adaptive per-parameter
scaling subsumes any Phi modulator's effect---provides a mechanistic
explanation supported by weight norm convergence analysis (1.2\%
variation under AdamW vs.~97\% under SGD), alignment informativeness
diagnostics (AIS invariant across methods), and the confirmed SGD
boundary condition.

Practitioners using AdamW can safely rely on constant weight decay---the
simplest strategy already captures the available benefit. Dynamic weight
decay research should prioritize conditions where Phi Invariance breaks:
SGD without batch normalization, ImageNet and LLM scale, and Vision
Transformer architectures where the conjecture's boundary conditions are
most likely violated.

\begin{center}\rule{0.5\linewidth}{0.5pt}\end{center}

\hypertarget{figures-and-tables}{%
\subsection{Figures and Tables}\label{figures-and-tables}}

\begin{itemize}
\tightlist
\item
  Figure 1: fig1\_taxonomy.png --- Phi Framework taxonomy showing
  methods along four modulation axes.
\item
  Figure 2: fig2\_accuracy\_comparison.png --- AdamW + ResNet-20
  accuracy comparison (CIFAR-10 and CIFAR-100).
\item
  Figure 3: multi\_arch\_comparison.png --- Multi-architecture accuracy
  comparison (ResNet-20 vs.~VGG-16-BN).
\item
  Figure 4: fig3\_bem\_vs\_accuracy.png --- BEM vs.~accuracy scatter
  plot.
\item
  Figure 5: fig4\_diagnostic\_heatmap.png --- CSI, AIS, BEM diagnostic
  heatmap.
\item
  Figure 6: fig6\_sgd\_vs\_adamw\_norms.png --- Weight norm comparison:
  AdamW vs.~SGD.
\item
  Figure 7: fig5\_weight\_norm\_trajectories.png --- Weight norm
  trajectories under AdamW.
\item
  Figure 8: certified\_band.png --- Certified band with method
  trajectories overlaid.
\item
  Figure 9: theorem2\_validation.png --- Cumulative vs.~worst-case
  alignment correlation.
\item
  Table 1: inline --- Method catalog (Phi modulator taxonomy).
\item
  Table 2: inline --- AdamW + ResNet-20 main accuracy results.
\item
  Table 3: inline --- Statistical tests vs.~constant baseline.
\item
  Table 4: inline --- SGD + ResNet-20 accuracy results.
\item
  Table 5: inline --- SGD + VGG-16-BN accuracy results.
\item
  Table 6: inline --- Diagnostic metrics (CSI, AIS, BEM).
\end{itemize}

\end{document}
